% =====================================================================
%  CHƯƠNG 4 — PHÂN TÍCH THIẾT KẾ, TRIỂN KHAI VÀ ĐÁNH GIÁ HỆ THỐNG
%  Bản nâng cấp: văn phong đặc tả kỹ thuật + bảng hoá + test cases.
%  Mọi số liệu kỹ thuật trích từ datn_agent_skills-test/project_setup/architecture/*
%  Các số liệu ĐO ĐẠC chưa có nguồn được đánh dấu  % TODO[cần đo]
% =====================================================================
\documentclass[../DoAn.tex]{subfiles}
\begin{document}

% Yêu cầu gói (đặt ở preamble DoAn.tex nếu chưa có):
%   \usepackage{booktabs}  % \toprule \midrule \bottomrule
%   \usepackage{array}     % cột p{} căn trên
%   \usepackage{float}     % [H]
%   \usepackage{longtable} % bảng test-case dài qua trang
%   \usepackage{makecell}  % xuống dòng trong ô

Chương này đặc tả quá trình phân tích, thiết kế kỹ thuật, triển khai và kiểm
thử hệ thống mô phỏng mạng truyền thông trên ô tô. Nội dung được trình bày
theo trình tự: (1) thiết kế kiến trúc — từ tổng thể đến phần cứng và phần mềm
kèm các lập luận lựa chọn giải pháp; (2) thiết kế chi tiết từng nút và đặc tả
giao thức truyền dữ liệu; (3) quy trình triển khai mã nguồn; và (4) các kịch
bản kiểm thử ở cấp độ module nhằm xác minh tính toàn vẹn và độ tin cậy của
toàn hệ thống.

\section{Thiết kế kiến trúc hệ thống}
\label{section:4.1}

Hệ thống được tổ chức theo mô hình phân tầng chức năng gồm ba nút độc lập,
mỗi nút giữ một trách nhiệm duy nhất và không chồng lấn: \textit{Nút cảm biến}
(Sensor Node) là nguồn sự thật về trạng thái xe, \textit{Nút cổng kết nối}
(Gateway Node) chỉ làm nhiệm vụ dịch giao thức, và \textit{Ứng dụng hiển thị}
(Android Cluster Application) chạy trên máy tính nhúng Raspberry Pi~4 đảm nhận
trực quan hoá. Việc tách bạch trách nhiệm này phản ánh đúng tô-pô ECU --
gateway -- head-unit của một mạng ô tô thực và cho phép kiểm thử từng nút một
cách độc lập.

% >>> HÌNH 4.1 — Sơ đồ khối kiến trúc tổng thể 3 tầng (Sensor → Gateway → Pi/Android)
% Nguồn: architecture/system_topology.md §1
\begin{figure}[H]
  \centering
  % \includegraphics[width=0.95\textwidth]{fig/architecture_overview}
  \caption{Sơ đồ khối kiến trúc tổng thể hệ thống ba tầng.}
  \label{fig:arch-overview}
\end{figure}

\subsection{Kiến trúc tổng thể}
\label{section:4.1.1}

Luồng dữ liệu của hệ thống là một chuỗi biến đổi một chiều, được tóm tắt trong
Bảng~\ref{tab:dataflow}. Nút cảm biến đọc đầu vào vật lý, sở hữu cấu trúc trạng
thái và phát hai khung CAN tiêu chuẩn (mã định danh \texttt{0x100} và
\texttt{0x200}). Nút cổng kết nối nhận khung CAN, tái dựng cấu trúc
\texttt{VehicleState\_t} 18~byte và chuyển tiếp nguyên khối qua TCP. Ứng dụng
Android, đóng vai trò máy chủ TCP, đọc đúng 18~byte mỗi gói, phân giải và vẽ
cụm đồng hồ số.

\begin{table}[H]
  \centering
  \caption{Luồng dữ liệu đầu cuối và ranh giới trách nhiệm (nguồn: integration.md).}
  \label{tab:dataflow}
  \begin{tabular}{@{}p{2.6cm}p{4.2cm}p{4.6cm}@{}}
    \toprule
    \textbf{Tầng} & \textbf{Trách nhiệm} & \textbf{Đầu ra / giao diện} \\
    \midrule
    Nút cảm biến (STM32-1st) & Đọc keypad + 3 biến trở (ADC), sở hữu
      \texttt{ControlData\_t}/\texttt{SensorData\_t} & CAN \texttt{0x100} (4\,B) +
      \texttt{0x200} (8\,B), 250\,kbps \\
    Nút cổng kết nối (STM32-2nd) & Kiểm tra + giải mã CAN, tích luỹ
      \texttt{VehicleState\_t}; \emph{không} chứa logic xe & TCP client gửi
      18\,byte thô tới \texttt{192.168.10.104:5000} \\
    Raspberry Pi 4 & Máy chủ chạy ứng dụng Android (host) & Cung cấp điểm cuối
      mạng \texttt{eth0} \\
    Ứng dụng Android & Phân giải gói, suy diễn cảnh báo hiển thị, vẽ giao diện &
      Cụm đồng hồ số $\sim$30\,FPS \\
    \bottomrule
  \end{tabular}
\end{table}

\subsubsection*{Lập luận lựa chọn kiến trúc}

\paragraph{Vì sao tách ba tầng độc lập.}
Mô hình một nút ghi -- nhiều nút đọc loại bỏ xung đột hợp nhất trạng thái: chỉ
Nút cảm biến được phép tạo ra giá trị, các nút phía sau là bộ chuyển tiếp/hiển
thị không trạng thái. Nhờ đó mỗi nút có thể được phát triển, nạp và kiểm thử
riêng rẽ, đồng thời tái hiện trung thực phân cấp ECU/gateway/head-unit của xe
thực -- vốn là mục tiêu của đề tài.

\paragraph{Vì sao dùng Gateway chuyển CAN sang Ethernet thay vì Wi-Fi.}
Liên kết giữa Nút cổng kết nối và máy tính nhúng được hiện thực bằng Ethernet
có dây điểm-điểm qua bộ điều khiển phần cứng W5500, thay vì Wi-Fi, dựa trên ba
luận điểm:
\begin{itemize}
  \item \textbf{Độ trễ tất định.} Đường truyền có dây loại bỏ tranh chấp môi
    trường vô tuyến, nhiễu RF và độ trễ dao động (jitter) -- yếu tố quyết định
    với một cụm đồng hồ thời gian thực.
  \item \textbf{Giải phóng tài nguyên vi điều khiển.} W5500 tích hợp sẵn ngăn
    xếp TCP/IP cứng; vi điều khiển chỉ thao tác qua API \texttt{socket / connect
    / send} mà không phải nhúng một ngăn xếp phần mềm (lwIP), tiết kiệm Flash/RAM
    vốn rất hạn chế trên STM32F103C8.
  \item \textbf{Độ tin cậy có sẵn.} TCP bảo đảm thứ tự và toàn vẹn gói tin nên
    không cần thêm cơ chế truyền lại ở lớp ứng dụng; trong khi Wi-Fi còn kéo theo
    chi phí kết hợp (association), bảo mật và chuyển vùng không cần thiết cho một
    liên kết cố định.
\end{itemize}

\paragraph{Vì sao dùng FreeRTOS để quản lý đa nhiệm.}
Trên cả hai nút STM32, hệ điều hành thời gian thực FreeRTOS (CMSIS-RTOS~v2) cho
phép phân tách các hoạt động có ràng buộc thời gian khác nhau thành các tác vụ
độc lập theo độ ưu tiên có quyền chiếm dụng (preemptive). Cơ chế hàng đợi
(queue) tách phần sinh dữ liệu khỏi phần truyền: ở Nút cảm biến, các tác vụ đọc
phím và đọc ADC chỉ nạp thông điệp vào một hàng đợi chung, còn một tác vụ duy
nhất chịu trách nhiệm truyền CAN -- nhờ đó tuần tự hoá truy cập bus và tránh
gọi đồng thời \texttt{HAL\_CAN\_AddTxMessage}. Ở Nút cổng kết nối, ngắt nhận
CAN chỉ đẩy khung vào hàng đợi an toàn-ngắt rồi trả quyền điều khiển ngay, để
tác vụ nền thực hiện phần xử lý nặng hơn -- đảm bảo độ phản hồi của ISR.

\subsection{Kiến trúc phần cứng}
\label{section:4.1.2}

Cả hai nút nhúng dùng chung vi điều khiển 32-bit STM32F103C8Tx (lõi ARM
Cortex-M3, tối đa 72\,MHz) tích hợp ngoại vi bxCAN. Bảng~\ref{tab:hw} liệt kê
cấu hình phần cứng đặc thù của từng nút.

\begin{table}[H]
  \centering
  \caption{Cấu hình phần cứng theo từng nút (nguồn: sensor\_ecu.md, gateway\_ecu.md, system\_topology.md).}
  \label{tab:hw}
  \begin{tabular}{@{}p{2.6cm}p{5.0cm}p{4.0cm}@{}}
    \toprule
    \textbf{Nút} & \textbf{Thành phần} & \textbf{Kết nối / tham số} \\
    \midrule
    \multirow Nút cảm biến
      & MCU STM32F103C8Tx & --- \\
      & Bộ thu phát CAN MCP2551 & CAN1: PA11 (RX) / PA12 (TX) \\
      & Ma trận phím 4$\times$3 & GPIO quét cột--hàng \\
      & 3 biến trở (chiết áp) & ADC1 IN0/IN1/IN2 qua DMA \\
      & 8 LED trạng thái & 74HC595 (dịch chốt) qua SPI \\
      & Gỡ lỗi & USART1 @ 115200 8N1 \\
    \midrule
    Nút cổng kết nối
      & MCU STM32F103C8Tx & --- \\
      & Bộ thu phát CAN MCP2551 & CAN1: PA11 (RX) / PA12 (TX) \\
      & Module Ethernet W5500 & SPI1, chip-select PB12 \\
      & Gỡ lỗi & USART1: PA9 (TX) / PA10 (RX) @ 115200 \\
    \midrule
    Head-unit
      & Raspberry Pi 4 + màn hình cảm ứng & Ethernet \texttt{eth0} nối thẳng W5500 \\
    \bottomrule
  \end{tabular}
\end{table}

Bộ ngoại vi CAN hoạt động ở mức logic TTL nên cần bộ thu phát MCP2551 chuyển
sang tín hiệu vi sai vật lý; hai điện trở đầu cuối 120\,$\Omega$ đặt tại hai đầu
mút bus để triệt phản xạ. Module W5500 kết nối qua SPI1 (chip-select PB12) và
tích hợp ngăn xếp TCP/IP cứng. Cổng \texttt{eth0} của Raspberry Pi~4 nối trực
tiếp tới W5500 bằng cáp mạng trên một mạng con riêng \texttt{192.168.10.0/24},
tách khỏi mạng Wi-Fi để tránh nhập nhằng định tuyến.

\medskip
\noindent\textbf{Ghi chú đính chính so với bản nháp:} bộ thu phát CAN là
\textbf{MCP2551} (không phải TJA1050) và các giá trị tốc độ/định thời CAN ở
\S\ref{section:4.3.1} dùng \textbf{250\,kbps} đúng theo cấu hình mã nguồn.

% TODO[cần đo]: tốc độ xung SPI thực tế của W5500 (bản nháp ghi 24 MHz nhưng
% chưa có trong tài liệu kiến trúc — xác nhận lại từ cấu hình hspi1).

\subsection{Kiến trúc phần mềm}
\label{section:4.1.3}

Phần mềm nhúng trên cả hai nút được tổ chức theo mô hình đa tác vụ FreeRTOS với
các hàng đợi và bộ định thời (timer) làm cơ chế đồng bộ. Hình~\ref{fig:rtos-task}
nên minh hoạ sơ đồ các tác vụ và hàng đợi của hai nút.

% >>> HÌNH 4.2 — Sơ đồ các Task/Queue trong FreeRTOS (Sensor + Gateway)
% Nguồn: sensor_ecu.md §2, gateway_ecu.md §2
\begin{figure}[H]
  \centering
  % \includegraphics[width=0.95\textwidth]{fig/rtos_tasks}
  \caption{Sơ đồ tác vụ và hàng đợi FreeRTOS trên Nút cảm biến và Nút cổng kết nối.}
  \label{fig:rtos-task}
\end{figure}

\section{Thiết kế chi tiết hệ thống}
\label{section:4.2}

\subsection{Thiết kế Nút cảm biến}
\label{section:4.2.1}

Nút cảm biến mô phỏng khối điều khiển khoang lái, thu thập hành vi người dùng và
duy trì trạng thái động học của xe. Phần mềm gồm bốn tác vụ FreeRTOS như
Bảng~\ref{tab:sensor-tasks}.

\begin{table}[H]
  \centering
  \caption{Các tác vụ FreeRTOS của Nút cảm biến (nguồn: sensor\_ecu.md §2).}
  \label{tab:sensor-tasks}
  \begin{tabular}{@{}p{3.0cm}p{2.0cm}p{2.4cm}p{4.4cm}@{}}
    \toprule
    \textbf{Tác vụ} & \textbf{Ưu tiên} & \textbf{Chu kỳ} & \textbf{Trách nhiệm} \\
    \midrule
    \texttt{buttonTask} & Low & $\sim$20\,ms & Quét keypad, khử rung, cập nhật
      \texttt{g\_ctrl\_data}, điều khiển LED, mã hoá \texttt{0x100} khi có phím \\
    \texttt{ADCTask} & BelowNormal & 100\,ms (10\,Hz) & Chép
      \texttt{adcBuffer[0..2]} $\to$ \texttt{g\_sensor\_data}, mã hoá
      \texttt{0x200}, đẩy hàng đợi \\
    \texttt{CANTask} & High & theo sự kiện & Lấy thông điệp khỏi hàng đợi, gọi
      \texttt{CAN\_Send} (\texttt{HAL\_CAN\_AddTxMessage}) \\
    \texttt{defaultTask} & Normal & theo sự kiện & Bộ phân tích lệnh UART
      (\texttt{PING}/\texttt{GET\_STATUS}/\dots) \\
    \bottomrule
  \end{tabular}
\end{table}

Hàng đợi chung \texttt{sensorQueue} (16 phần tử \texttt{CAN\_Message\_t}) nhận
thông điệp từ hai nguồn (\texttt{ADCTask} phát \texttt{0x200},
\texttt{buttonTask} phát \texttt{0x100}) và được tiêu thụ bởi một mình
\texttt{CANTask}, qua đó tuần tự hoá toàn bộ truyền CAN.

\medskip
\noindent\textbf{Phân định nguồn tín hiệu (đính chính bản nháp):} ba đại lượng
tương tự -- vận tốc, mức nhiên liệu và vòng tua máy -- được lấy từ \textbf{ba
biến trở qua ADC/DMA} (\texttt{adcBuffer[0]=speed}, \texttt{[1]=fuel},
\texttt{[2]=rpm}), \emph{không} suy ra từ thời gian giữ phím. Ma trận phím chỉ
\emph{bật/tắt} (toggle) các trường điều khiển: đèn xi-nhan, đèn pha/cốt, hộp số,
chế độ lái, cửa, dây an toàn, sấy kính, v.v. (Bảng~\ref{tab:keypad}).

\begin{table}[H]
  \centering
  \caption{Ánh xạ phím $\to$ hành động trên Nút cảm biến (nguồn: sensor\_ecu.md §3).}
  \label{tab:keypad}
  \begin{tabular}{@{}cll@{}}
    \toprule
    \textbf{Phím} & \textbf{Hành động} & \textbf{Khoá liên động} \\
    \midrule
    1 & turnL & xoá turnR, hazard \\
    2 & hazard & xoá turnL, turnR \\
    3 & turnR & xoá turnL, hazard \\
    4 & highBeam & xoá lowBeam \\
    5 & gear (P$\to$R$\to$N$\to$D) & --- \\
    6 & lowBeam & xoá highBeam \\
    7 & seatbelt & --- \\
    8 & mode (lặp vòng) & --- \\
    9 & door & --- \\
    * & heat & --- \\
    0 & steer & --- \\
    \# & wind & --- \\
    \bottomrule
  \end{tabular}
\end{table}

Để khử rung tiếp điểm cơ khí, một phím chỉ được ghi nhận khi trạng thái đồng
nhất qua hai chu kỳ quét liên tiếp. Các kênh ADC dùng thời gian lấy mẫu
\texttt{ADC\_SAMPLETIME\_239CYCLES\_5} ($\approx$21\,$\mu$s/kênh @ 12\,MHz) nhằm
tránh hiện tượng rò điện tích giữa các kênh khi nguồn biến trở có trở kháng cao.

\subsection{Thiết kế Nút cổng kết nối}
\label{section:4.2.2}

Nút cổng kết nối là cầu nối định tuyến dữ liệu bất đối xứng giữa mạng CAN hướng
gói tin và mạng Ethernet/TCP hướng dòng byte. Lõi xử lý dựa trên hàng đợi
\texttt{canToEthQueue} (16 phần tử) làm bộ đệm chống lệch pha tốc độ:

\begin{enumerate}
  \item Khung CAN hợp lệ vượt bộ lọc phần cứng $\to$ ngắt \texttt{USB\_LP\_CAN1\_RX0\_IRQn}
    (NVIC ưu tiên 5) đẩy khung vào \texttt{canToEthQueue} rồi trả quyền điều khiển ngay.
  \item \texttt{Eth\_SendTask} lấy khung khỏi hàng đợi, gọi
    \texttt{CAN\_ValidateFrame} (ID\,$\le$\,0x7FF, DLC\,$\le$\,8).
  \item \texttt{decodeCAN} kiểm CRC8, giải bit cho \texttt{0x100} hoặc giải
    big-endian cho \texttt{0x200}, theo dõi số thứ tự (sequence) để phát hiện
    mất/trùng khung.
  \item Tích luỹ lai: một \texttt{VehicleState\_t} bền vững giữ \emph{control mới
    nhất} và \emph{sensor mới nhất}; mỗi khung hợp lệ kích hoạt gửi nguyên ảnh
    chụp 18~byte.
\end{enumerate}

Máy trạng thái socket TCP-client (không chặn, \texttt{SF\_IO\_NONBLOCK}) trên
socket~0 vận hành theo Bảng~\ref{tab:sock-fsm}. Chế độ không chặn bảo đảm
\texttt{connect()} không khoá \texttt{Eth\_SendTask}, nhờ đó hàng đợi không tràn
khi ứng dụng chưa sẵn sàng.

\begin{table}[H]
  \centering
  \caption{Máy trạng thái socket của Nút cổng kết nối (nguồn: gateway\_ecu.md §4).}
  \label{tab:sock-fsm}
  \begin{tabular}{@{}lll@{}}
    \toprule
    \textbf{Trạng thái (\texttt{getSn\_SR})} & \textbf{Hành động} & \textbf{Trả về} \\
    \midrule
    \texttt{SOCK\_ESTABLISHED} & sẵn sàng gửi & 1 \\
    \texttt{SOCK\_CLOSED} & \texttt{socket(0,TCP,port++,SF\_IO\_NONBLOCK)} & 0 \\
    \texttt{SOCK\_INIT} & \texttt{connect(0, APP\_IP, APP\_PORT)} & 0 \\
    \texttt{SOCK\_SYNSENT} & chờ bắt tay, không gọi lại \texttt{connect} & 0 \\
    \texttt{SOCK\_CLOSE\_WAIT} & \texttt{disconnect(0)} & 0 \\
    \bottomrule
  \end{tabular}
\end{table}

% >>> HÌNH 4.3 — Lưu đồ pipeline RX→decode→accumulate→TCP của Gateway
% Nguồn: gateway_ecu.md §3
\begin{figure}[H]
  \centering
  % \includegraphics[width=0.85\textwidth]{fig/gateway_pipeline}
  \caption{Lưu đồ xử lý của Nút cổng kết nối: nhận CAN, giải mã, tích luỹ và gửi TCP.}
  \label{fig:gw-pipeline}
\end{figure}

\subsection{Thiết kế ứng dụng Android}
\label{section:4.2.3}

Ứng dụng đóng vai trò cụm đồng hồ số và \emph{máy chủ} TCP. Khác với bản nháp,
ứng dụng là phía \textbf{server} (\texttt{ServerSocket} trên
\texttt{0.0.0.0:5000}), còn W5500 của Nút cổng kết nối là phía \textbf{client}
chủ động kết nối ra; head-unit là điểm lắng nghe ổn định, thiết bị nhúng tự kết
nối lại mỗi 2\,giây.

\texttt{GatewayClient} là singleton phạm vi tiến trình; vòng đời máy chủ thuộc
\texttt{SocketService} (foreground service, \texttt{START\_STICKY}) để cổng 5000
chỉ bind một lần và sống sót khi hệ điều hành thu hồi tiến trình. Luồng nhận
\texttt{readFully(18)} gom đúng một gói 18~byte (xử lý phân mảnh dòng TCP) rồi
gọi \texttt{applyGatewayPacket}. Sau khi phân giải, mô hình \texttt{VehicleState}
(POJO) được đồng bộ sang luồng giao diện và vẽ bằng \texttt{GaugeView}.

Ngoài các trường nhận từ dây, ứng dụng tự \emph{suy diễn} một số cảnh báo hiển
thị (Bảng~\ref{tab:derived}); các cảnh báo này không do firmware gửi, nhằm giữ
tải tin tối thiểu và cho phép tinh chỉnh ngưỡng ở tầng giao diện.

\begin{table}[H]
  \centering
  \caption{Cảnh báo suy diễn tại ứng dụng (nguồn: android\_app.md §3).}
  \label{tab:derived}
  \begin{tabular}{@{}ll@{}}
    \toprule
    \textbf{Cảnh báo} & \textbf{Quy tắc} \\
    \midrule
    \texttt{engineCheckOn} & rpm $>$ 6500 \\
    \texttt{breakOn} & giảm tốc $>$ 12 km/h mỗi giây giữa hai gói \\
    \texttt{oilPressureOn} & rpm $<$ 500 \emph{và} speed $>$ 5 \\
    \texttt{temperatureOn} & mô hình nhiệt ảo vượt ngưỡng 0.85 \\
    \bottomrule
  \end{tabular}
\end{table}

\subsection{Thiết kế giao thức truyền dữ liệu}
\label{section:4.2.4}

Hệ thống định nghĩa hai tầng giao thức tách biệt: tầng CAN (giữa Nút cảm biến và
Nút cổng kết nối) và tầng TCP/Ethernet (giữa Nút cổng kết nối và ứng dụng).

\subsubsection{Tầng CAN}

Tầng CAN dùng khung chuẩn 11-bit, 250\,kbps. Cấu trúc định danh phản ánh độ ưu
tiên: khung động học \texttt{0x200} (cao) truyền định kỳ 10\,Hz, khung điều
khiển \texttt{0x100} truyền theo sự kiện khi có phím. Mỗi khung mang một
\textbf{CRC8} (đa thức 0x07, MSB-first, init 0x00) và một bộ đếm \textbf{số thứ
tự 4-bit} để phát hiện hỏng/mất/trùng khung ở lớp ứng dụng.
Bảng~\ref{tab:can-0x100} và Bảng~\ref{tab:can-0x200} đặc tả bố cục byte.

\begin{table}[H]
  \centering
  \caption{Khung \texttt{0x100} -- CONTROL, DLC 4 (nguồn: can\_database.md).}
  \label{tab:can-0x100}
  \begin{tabular}{@{}ccl@{}}
    \toprule
    \textbf{Byte} & \textbf{Bit} & \textbf{Trường} \\
    \midrule
    0 & 0--7 & turnL, turnR, hazard, highBeam, lowBeam, door, seatbelt, heat \\
    1 & 0--1 & gear (0=P,1=R,2=N,3=D) \\
    1 & 2--3 & mode (0=ECO,1=SPORT,2=COMFORT,3=NORMAL) \\
    1 & 4--5 & steer, wind \\
    2 & 0--3 & sequence (nibble thấp) \\
    3 & --- & CRC8 trên byte 0--2 \\
    \bottomrule
  \end{tabular}
\end{table}

\begin{table}[H]
  \centering
  \caption{Khung \texttt{0x200} -- SENSOR, DLC 8, giá trị 16-bit big-endian (nguồn: can\_database.md).}
  \label{tab:can-0x200}
  \begin{tabular}{@{}cl@{}}
    \toprule
    \textbf{Byte} & \textbf{Trường} \\
    \midrule
    0--1 & speed (MSB, LSB) \\
    2--3 & fuel (MSB, LSB) \\
    4--5 & rpm (MSB, LSB) \\
    6 & sequence (nibble cao) \\
    7 & CRC8 trên byte 0--6 \\
    \bottomrule
  \end{tabular}
\end{table}

\subsubsection{Tầng TCP/Ethernet}

Tải tin Ethernet là một cấu trúc \texttt{VehicleState\_t} cố định 18~byte. Nút
cổng kết nối \texttt{memcpy} nguyên cấu trúc này vào bộ đệm rồi gọi
\texttt{send()}; phía Android dùng \texttt{readFully(18)} để gom đúng một gói.
Do TCP đã bảo đảm thứ tự và toàn vẹn của dòng byte, hệ thống không bổ sung byte
đồng bộ, mã loại thông điệp hay byte kiểm tra ở lớp ứng dụng; toàn vẹn được giao
phó cho TCP. Kích thước cố định cho phép phân giải bằng cơ chế length-framing
theo độ dài định trước. Bố cục 18~byte little-endian được đặc tả trong
Bảng~\ref{tab:tcp-packet}.

\begin{table}[H]
  \centering
  \caption{Cấu trúc gói TCP 18 byte (\texttt{VehicleState\_t}, little-endian) -- nguồn: vehicle\_state.md §2.}
  \label{tab:tcp-packet}
  \begin{tabular}{@{}cllp{5.2cm}@{}}
    \toprule
    \textbf{Byte} & \textbf{Trường} & \textbf{Kiểu} & \textbf{Ý nghĩa / biến đổi tại Android} \\
    \midrule
    0  & turnL    & uint8 & đèn xi-nhan trái ($\ne 0$) \\
    1  & turnR    & uint8 & đèn xi-nhan phải ($\ne 0$) \\
    2  & hazard   & uint8 & đèn khẩn cấp \\
    3  & highBeam & uint8 & đèn pha \\
    4  & lowBeam  & uint8 & đèn cốt \\
    5  & door     & uint8 & cảnh báo cửa mở \\
    6  & seatbelt & uint8 & cảnh báo dây an toàn \\
    7  & gear     & uint8 & số (P/R/N/D) \\
    8  & mode     & uint8 & chế độ lái (ECO/SPORT/COMFORT/NORMAL) \\
    9  & heat     & uint8 & sấy kính (defrost) \\
    10 & steer    & uint8 & sấy vô-lăng \\
    11 & wind     & uint8 & sấy kính sau \\
    12--13 & speed & uint16 LE & vận tốc thô; Android: $\dfrac{\text{giá trị}}{4095}\times240$ km/h \\
    14--15 & fuel  & uint16 LE & nhiên liệu thô; Android: $\dfrac{\text{giá trị}}{4095}$ (0--1) \\
    16--17 & rpm   & uint16 LE & vòng tua thô; Android: $\dfrac{\text{giá trị}}{4095}\times8000$ \\
    \bottomrule
  \end{tabular}
\end{table}

\noindent\textbf{Lưu ý endianness:} khung CAN \texttt{0x200} mang speed/fuel/rpm
\emph{big-endian}; Nút cổng kết nối tái dựng về \texttt{uint16} bản địa
($\texttt{data[0]}\ll8 \mid \texttt{data[1]}$); \texttt{memcpy} cấu trúc phát ra
\emph{little-endian} (bản địa Cortex-M3) và Android cũng đọc little-endian. Giá
trị được bảo toàn đầu cuối; đây là một \emph{cặp} biến đổi khớp nhau, không được
sửa lệch một phía.

\subsubsection*{Hệ số chuẩn hoá thang đo}

Cảm biến gửi \emph{thẳng} giá trị ADC 12-bit (0--4095), không phóng lên 16-bit;
do đó ứng dụng chia cho toàn thang \textbf{4095} (hằng số
\texttt{ClusterController.ADC\_FULL\_SCALE}). Vận tốc cực đại biểu diễn được do
đó là 240\,km/h.

\section{Triển khai hệ thống}
\label{section:4.3}

\subsection{Triển khai Nút cảm biến}
\label{section:4.3.1}

Mã nguồn được cấu hình bằng STM32CubeMX và biên dịch bằng STM32CubeIDE trên nền
thư viện HAL. Nhân FreeRTOS được khởi tạo, cấp phát bộ nhớ theo mô hình tĩnh kết
hợp động. Bộ ngoại vi bxCAN được cấu hình để đạt tốc độ \textbf{250\,kbps}: với
bus ngoại vi APB1 ở 36\,MHz, đặt \textbf{bộ chia (prescaler) = 9}, phân đoạn
\textbf{BS1 = 13\,TQ}, \textbf{BS2 = 2\,TQ}, \textbf{SJW = 1\,TQ}, tổng
\textbf{16\,TQ/bit} $\Rightarrow$ 4\,$\mu$s/bit = 250\,kbps. Việc truyền được
thực hiện bằng cách điền cấu trúc tiêu đề truyền, chọn mailbox trống và gọi
\texttt{HAL\_CAN\_AddTxMessage}. Cả hai nút bật \texttt{AutoBusOff} và
\texttt{AutoRetransmission} để tự phục hồi trên liên kết telemetry hai nút.

\medskip
\noindent\textbf{Đính chính bản nháp:} thông số định thời CAN là
prescaler~9 / BS1~13 / BS2~2 (16~TQ, 250\,kbps), thay cho prescaler~4 / BS1~13 /
BS2~4 (18~TQ, 500\,kbps) trong bản nháp.

\subsection{Triển khai Nút cổng kết nối}
\label{section:4.3.2}

Triển khai tập trung vào cấu hình bộ lọc CAN và lớp vật lý cho W5500. Bộ lọc CAN
được đặt ở chế độ danh sách định danh, chỉ cho các khung \texttt{0x100}/\texttt{0x200}
đi vào FIFO nhằm giảm tần suất ngắt. W5500 được khởi tạo với các callback SPI
(đọc/ghi byte, chọn chip PB12) và cấu hình IP tĩnh
(\texttt{192.168.10.50}, gateway \texttt{192.168.10.1}, mask
\texttt{255.255.255.0}, \texttt{NETINFO\_STATIC}). Socket được mở ở chế độ TCP
client không chặn, \texttt{connect} tới \texttt{192.168.10.104:5000}
(\texttt{eth0} của Pi); sau 5 lần gửi lỗi liên tiếp, socket được đóng để mở lại
sạch sẽ.

\medskip
\noindent\textbf{Lưu ý thứ tự khởi tạo:} \texttt{HAL\_CAN\_Start} và
\texttt{HAL\_CAN\_ActivateNotification} được gọi ở đầu \texttt{Eth\_SendTask}
(sau khi hàng đợi đã tạo và scheduler chạy), không gọi trong \texttt{main()}
trước scheduler, nhằm tránh ngắt RX đẩy vào hàng đợi NULL gây treo
\texttt{configASSERT}.

\subsection{Triển khai ứng dụng Android}
\label{section:4.3.3}

Ứng dụng được viết bằng Java (Android SDK min~24 / target~36, Gradle). Một dịch
vụ nền (\texttt{SocketService}, foreground, \texttt{START\_STICKY}) duy trì máy
chủ TCP; luồng nhận \texttt{TCP\_Recv\_Thread} chạy ở \texttt{MAX\_PRIORITY},
dùng \texttt{DataInputStream.readFully} với bộ đệm cố định để tránh cấp phát lặp
gây kích hoạt bộ dọn rác. \texttt{BootReceiver} tự khởi động khi máy bật,
\texttt{WatchdogReceiver} hồi sinh dịch vụ chết. Dữ liệu phân giải xong được
chuyển an toàn sang luồng giao diện và kích hoạt vẽ lại \texttt{GaugeView}.

\subsection{Thư viện và công cụ sử dụng}
\label{section:4.3.4}

Bảng~\ref{tab:tools} tổng hợp công cụ và thư viện sử dụng trong toàn dự án.

\begin{table}[H]
  \centering
  \caption{Công cụ và thư viện sử dụng (nguồn: overview.md §3, các doc kiến trúc).}
  \label{tab:tools}
  \begin{tabular}{@{}p{3.6cm}p{2.6cm}p{6.0cm}@{}}
    \toprule
    \textbf{Công cụ / thư viện} & \textbf{Phạm vi} & \textbf{Vai trò} \\
    \midrule
    STM32CubeMX & 2 nút STM32 & Cấu hình đồ hoạ ngoại vi, sinh mã khởi tạo HAL \\
    STM32CubeIDE (GCC ARM) & 2 nút STM32 & Biên dịch, nạp và gỡ lỗi qua ST-Link \\
    STM32 HAL & 2 nút STM32 & Lớp trừu tượng hoá phần cứng (GPIO/SPI/ADC/CAN/UART) \\
    FreeRTOS (CMSIS-RTOS v2) & 2 nút STM32 & Đa nhiệm, hàng đợi, timer, mutex \\
    WIZnet ioLibrary & Gateway & Driver W5500, API socket TCP/IP cứng \\
    Android Studio + Gradle & ClusterApp & Phát triển, build APK \\
    Android SDK (Java) & ClusterApp & Dịch vụ nền, socket, giao diện \texttt{Canvas} \\
    Wireshark & Kiểm thử & Bắt và phân tích gói TCP trên LAN \\
    Máy hiện sóng / phân tích logic & Kiểm thử & Đo tín hiệu vi sai CAN, giải mã khung \\
    \bottomrule
  \end{tabular}
\end{table}

\section{Kiểm thử trong quá trình phát triển}
\label{section:4.4}

Công tác kiểm thử được tổ chức theo chiến lược \textit{kiểm thử từ dưới lên}
(bottom-up): trước hết xác minh từng tầng truyền thông một cách độc lập (CAN,
rồi Gateway CAN--Ethernet, rồi ứng dụng hiển thị), sau cùng kiểm thử tích hợp
toàn tuyến đầu cuối. Mỗi hạng mục được mô tả dưới dạng kịch bản kiểm thử (test
case) có mã định danh, điều kiện tiên quyết, các bước thực hiện và kết quả mong
đợi nhằm bảo đảm khả năng tái lập. Các giá trị đo thực nghiệm chưa chốt được
đánh dấu để tác giả điền sau khi đo.

\subsubsection*{Môi trường và công cụ kiểm thử}

Bảng~\ref{tab:test-env} liệt kê cấu hình môi trường dùng cho toàn bộ các kịch
bản kiểm thử trong mục này.

\begin{table}[H]
  \centering
  \caption{Môi trường và công cụ kiểm thử.}
  \label{tab:test-env}
  \begin{tabular}{@{}p{3.4cm}p{8.6cm}@{}}
    \toprule
    \textbf{Hạng mục} & \textbf{Cấu hình} \\
    \midrule
    Thiết bị thử & Nút cảm biến, Nút cổng kết nối (STM32F103C8), Raspberry Pi 4 + màn hình cảm ứng \\
    Mạng CAN & Bus 250\,kbps, MCP2551, 2 điện trở đầu cuối 120\,$\Omega$ \\
    Mạng Ethernet & W5500 nối thẳng \texttt{eth0} của Pi, subnet \texttt{192.168.10.0/24} \\
    Đo tín hiệu CAN & Máy hiện sóng số + bộ phân tích logic có giải mã CAN \\
    Phân tích gói TCP & Wireshark trên máy tính cùng LAN \\
    Quan sát phía nhúng & UART gỡ lỗi (115200 8N1): nhãn \texttt{[CAN RX]}, \texttt{[DIAG]}, \texttt{[ETH]} \\
    Quan sát phía Android & \texttt{DebugActivity} (tốc độ gói, bộ đếm lỗi), \texttt{adb logcat} \\
    \bottomrule
  \end{tabular}
\end{table}

\subsection{Kiểm thử truyền thông CAN}
\label{section:4.4.1}

\textbf{Mục tiêu.} Xác minh tính toàn vẹn của tín hiệu điện vật lý trên bus và
tính chính xác của cấu trúc khung \texttt{0x100}/\texttt{0x200}, bao gồm cơ chế
kiểm lỗi CRC8 và theo dõi số thứ tự ở lớp ứng dụng. \textbf{Phương pháp.} Dùng
máy hiện sóng quan sát mức vi sai CANH/CANL, bộ phân tích logic giải mã khung,
và UART của Gateway đọc bộ chẩn đoán \texttt{canDecodeDiag}. Các kịch bản chi
tiết trong Bảng~\ref{tab:tc-can}.

\begin{table}[H]
  \centering
  \caption{Kịch bản kiểm thử tầng CAN.}
  \label{tab:tc-can}
  \begin{tabular}{@{}p{1.7cm}p{2.6cm}p{4.2cm}p{4.2cm}@{}}
    \toprule
    \textbf{ID} & \textbf{Điều kiện tiên quyết} & \textbf{Các bước} & \textbf{Kết quả mong đợi} \\
    \midrule
    TC\_CAN\_01 & Hai nút cấp nguồn, bus có 2 điện trở 120\,$\Omega$; nối hiện
      sóng vào CANH/CANL & 1. Cho Nút cảm biến phát khung. 2. Quan sát mức điện
      áp vi sai. & Trạng thái lặn: CANH$\approx$CANL$\approx$2.5\,V (vi sai
      $\approx$0\,V). Trạng thái trội: CANH$\approx$3.5\,V, CANL$\approx$1.5\,V
      (vi sai $\approx$2.0\,V). \\
    TC\_CAN\_02 & Bộ phân tích logic + phần mềm giải mã CAN nối song song bus &
      1. Phát \texttt{0x200} định kỳ. 2. Giải mã khung. & Bit time đúng 4\,$\mu$s
      (250\,kbps); ID, DLC, byte tải khớp dữ liệu nguồn. % TODO[cần đo]: sai số bit-timing thực
      \\
    TC\_CAN\_03 & Kết nối UART gỡ lỗi của Gateway & 1. Phát liên tục
      \texttt{0x100}/\texttt{0x200}. 2. Đọc \texttt{[DIAG]}. & CRC8 khớp,
      \texttt{crc\_error}=0, không \texttt{seq\_gap}/\texttt{seq\_duplicate}. \\
    TC\_CAN\_04 & Như TC\_CAN\_03 & 1. Rút/chập một nhánh bus gây lỗi khung. &
      Khung lỗi bị loại; \texttt{crc\_error} hoặc \texttt{seq\_gap} tăng; bus tự
      phục hồi (AutoBusOff). \\
    \bottomrule
  \end{tabular}
\end{table}

\textbf{Nhận xét.} Mức vi sai $\approx$2.0\,V ở trạng thái trội và $\approx$0\,V
ở trạng thái lặn (TC\_CAN\_01) phù hợp với chuẩn CAN ô tô. Bộ đếm
\texttt{crc\_error}/\texttt{seq\_gap}/\texttt{seq\_duplicate} bằng 0 trong điều
kiện bình thường (TC\_CAN\_03) cho thấy đường truyền và thuật toán mã hoá/giải
mã hoạt động ổn định; khi chủ động gây lỗi (TC\_CAN\_04), khung hỏng bị loại và
bus tự phục hồi nhờ \texttt{AutoBusOff}.
% TODO[cần đo]: chèn ảnh chụp dạng sóng oscilloscope + ảnh giải mã từ logic analyzer.

\subsection{Kiểm thử Gateway CAN--Ethernet}
\label{section:4.4.2}

\textbf{Mục tiêu.} Xác minh năng lực dịch giao thức thời gian thực của Nút cổng
kết nối và khả năng chịu tải của hàng đợi phần mềm khi lưu lượng CAN tăng cao,
đồng thời kiểm tra cơ chế tự kết nối lại của máy trạng thái socket.
\textbf{Phương pháp.} Dùng Wireshark bắt gói TCP do W5500 phát, đối chiếu tải
tin 18~byte với dữ liệu nguồn; tạo kịch bản tải cao bằng cách hạ chu kỳ phát CAN;
và quan sát log \texttt{[ETH]} trên UART để theo dõi chuyển trạng thái socket.
Các kịch bản chi tiết trong Bảng~\ref{tab:tc-eth}.

\begin{table}[H]
  \centering
  \caption{Kịch bản kiểm thử Nút cổng kết nối (CAN $\to$ TCP).}
  \label{tab:tc-eth}
  \begin{tabular}{@{}p{1.7cm}p{2.6cm}p{4.2cm}p{4.2cm}@{}}
    \toprule
    \textbf{ID} & \textbf{Điều kiện tiên quyết} & \textbf{Các bước} & \textbf{Kết quả mong đợi} \\
    \midrule
    TC\_ETH\_01 & PC chạy Wireshark cùng LAN; W5500 đã \texttt{ESTABLISHED} &
      1. Bắt gói TCP từ Gateway. 2. Xét tải tin. & Mỗi gói đúng \textbf{18 byte};
      byte 0--11 là cờ điều khiển/số/chế độ, byte 12--17 là speed/fuel/rpm
      little-endian khớp nguồn. \emph{(Không có byte đồng bộ/CRC -- đúng thiết
      kế raw-struct.)} \\
    TC\_ETH\_02 & Như TC\_ETH\_01 & 1. Hạ chu kỳ phát CAN để tạo tải cao
      (stress). 2. Theo dõi liên tục. & Không tràn \texttt{canToEthQueue};
      không mất gói TCP trên LAN nội bộ. % TODO[cần đo]: throughput / tỉ lệ mất gói
      \\
    TC\_ETH\_03 & Ứng dụng Android chưa mở cổng :5000 & 1. Khởi động Gateway
      trước. 2. Đọc UART \texttt{[ETH]}. & Socket lặp \texttt{CLOSED}$\to$\texttt{INIT}
      $\to$\texttt{SYNSENT}; \texttt{Eth\_SendTask} không bị chặn (non-block). \\
    TC\_ETH\_04 & W5500 nối \texttt{eth0} Pi, subnet \texttt{192.168.10.0/24} &
      1. Mở ứng dụng. 2. Quan sát chuyển \texttt{ESTABLISHED}. & Log
      \texttt{[ETH] socket ESTABLISHED}; bắt đầu gửi 18~byte đều đặn. \\
    TC\_ETH\_05 & Đang \texttt{ESTABLISHED} & 1. Ngắt cáp 5\,s rồi cắm lại. &
      Sau $\le$5 lần gửi lỗi, socket đóng và mở lại; kết nối tự phục hồi. \\
    \bottomrule
  \end{tabular}
\end{table}

\medskip
\noindent\textbf{Đính chính bản nháp:} kết quả Wireshark phải cho thấy gói
\textbf{18 byte raw struct}, không phải khung có 2 byte đồng bộ + mã loại + byte
XOR. Mô tả "byte đồng bộ" và "checksum XOR" trong bản nháp thuộc đặc tả cũ đã bị
thay thế (xem D-010).

\textbf{Nhận xét.} Khi đối chiếu trên Wireshark (TC\_ETH\_01), mỗi gói TCP đúng
18~byte với trường speed/fuel/rpm little-endian trùng khớp dữ liệu gốc ở Nút cảm
biến. Trong kịch bản tải cao (TC\_ETH\_02), hàng đợi \texttt{canToEthQueue}
không tràn nhờ cơ chế socket không chặn (\texttt{SF\_IO\_NONBLOCK}) giữ cho
\texttt{Eth\_SendTask} luôn rút được hàng đợi. Các kịch bản TC\_ETH\_03 đến
TC\_ETH\_05 xác nhận máy trạng thái socket tự tiến tới \texttt{ESTABLISHED} và tự
phục hồi sau khi mất kết nối.

\subsection{Kiểm thử ứng dụng Android}
\label{section:4.4.3}

\textbf{Mục tiêu.} Bảo đảm tính ổn định của kết nối socket, độ chính xác của
thuật toán phân giải dòng byte và hiệu năng hiển thị thời gian thực của giao
diện. \textbf{Phương pháp.} Chạy ứng dụng trên Raspberry Pi 4, dùng
\texttt{DebugActivity}/\texttt{adb logcat} giám sát tốc độ gói và bộ đếm lỗi;
kiểm thử chéo bằng cách thao tác biến trở và phím trên Nút cảm biến rồi quan sát
giao diện. Các kịch bản chi tiết trong Bảng~\ref{tab:tc-app}.

\begin{table}[H]
  \centering
  \caption{Kịch bản kiểm thử ứng dụng hiển thị.}
  \label{tab:tc-app}
  \begin{tabular}{@{}p{1.7cm}p{2.6cm}p{4.2cm}p{4.2cm}@{}}
    \toprule
    \textbf{ID} & \textbf{Điều kiện tiên quyết} & \textbf{Các bước} & \textbf{Kết quả mong đợi} \\
    \midrule
    TC\_APP\_01 & Ứng dụng chạy trên Pi 4 + màn cảm ứng; Gateway phát dữ liệu &
      1. Mở ứng dụng. 2. Theo dõi \texttt{DebugActivity}. & Ghi nhận kết nối
      thành công; luồng nhận đếm gói đều; tốc độ gói ổn định. % TODO[cần đo]: packet rate thực
      \\
    TC\_APP\_02 & Như TC\_APP\_01 & 1. Vặn biến trở "speed" toàn dải. & Kim đồng
      hồ tốc độ quét mượt 0--240\,km/h; giá trị số khớp công thức $/4095\times240$. \\
    TC\_APP\_03 & Như TC\_APP\_01 & 1. Bấm các phím đèn/số/chế độ trên Nút cảm
      biến. & Biểu tượng tương ứng và nhãn số/chế độ cập nhật đúng; đồng bộ với
      LED tại chỗ. \\
    TC\_APP\_04 & Như TC\_APP\_01 & 1. Đo độ trễ từ thao tác vật lý đến thay đổi
      trên màn hình. & Độ trễ ở mức cảm nhận thời gian thực. % TODO[cần đo]: độ trễ end-to-end (ms)
      \\
    TC\_APP\_05 & Gửi gói lỗi/đứt giữa chừng & 1. Ngắt kết nối khi đang chạy. &
      Ứng dụng không treo; tự kết nối lại sau $\sim$2\,s (\texttt{WatchdogReceiver}). \\
    \bottomrule
  \end{tabular}
\end{table}

\textbf{Nhận xét.} Khi vặn biến trở (TC\_APP\_02), kim tốc độ quét mượt trong
dải 0--240\,km/h và giá trị số khớp công thức chuẩn hoá $/4095\times240$, xác
nhận hệ số thang đo đúng. Các thao tác đèn/số/chế độ (TC\_APP\_03) phản ánh tức
thời trên biểu tượng và nhãn, đồng bộ với LED tại chỗ ở Nút cảm biến.

% >>> HÌNH 4.4 — Ảnh chụp giao diện cụm đồng hồ khi chạy (kèm chú thích các vùng)
\begin{figure}[H]
  \centering
  % \includegraphics[width=0.9\textwidth]{fig/cluster_ui}
  \caption{Giao diện cụm đồng hồ số khi hệ thống vận hành.}
  \label{fig:cluster-ui}
\end{figure}

\subsection{Kiểm thử tích hợp đầu cuối}
\label{section:4.4.4}

\textbf{Mục tiêu.} Sau khi từng tầng đã được xác minh độc lập, mục này kiểm thử
toàn tuyến \textit{thao tác vật lý $\to$ CAN $\to$ Gateway $\to$ TCP $\to$ hiển
thị} như một hệ thống hoàn chỉnh, tập trung vào tính đúng đắn đầu cuối, độ trễ
cảm nhận và độ bền vững khi gặp sự cố mạng. Các kịch bản trong
Bảng~\ref{tab:tc-e2e} cũng bao quát những điểm dễ hỏng khi tích hợp đã ghi nhận
trong quá trình phát triển.

\begin{table}[H]
  \centering
  \caption{Kịch bản kiểm thử tích hợp đầu cuối.}
  \label{tab:tc-e2e}
  \begin{tabular}{@{}p{1.7cm}p{2.6cm}p{4.2cm}p{4.2cm}@{}}
    \toprule
    \textbf{ID} & \textbf{Điều kiện tiên quyết} & \textbf{Các bước} & \textbf{Kết quả mong đợi} \\
    \midrule
    TC\_E2E\_01 & Cả ba nút hoạt động, kết nối \texttt{ESTABLISHED} & 1. Vặn biến
      trở và bấm phím đồng thời. 2. Quan sát cụm đồng hồ. & Mọi đại lượng và biểu
      tượng cập nhật nhất quán; không lệch trường nào. \\
    TC\_E2E\_02 & Như TC\_E2E\_01 & 1. Đo thời gian từ thao tác đến thay đổi pixel
      trên màn hình. & Độ trễ đầu cuối ở mức cảm nhận thời gian thực.
      % TODO[cần đo]: độ trễ end-to-end (ms)
      \\
    TC\_E2E\_03 & Đang chạy ổn định & 1. Cấu hình sai \texttt{APP\_IP} (IP WiFi
      thay vì \texttt{eth0}). & W5500 kẹt \texttt{SYNSENT}, UART báo
      \texttt{[ETH] socket not ready}; giao diện chỉ hiện 3 gauge — đúng hành vi
      khi chưa có dữ liệu. \\
    TC\_E2E\_04 & Đang \texttt{ESTABLISHED} & 1. Khởi động lại ứng dụng Android
      trong khi Gateway vẫn phát. & Gateway tự kết nối lại; ứng dụng tiếp tục
      nhận sau khi cổng \texttt{:5000} mở lại (\texttt{START\_STICKY}). \\
    TC\_E2E\_05 & Đang chạy ổn định & 1. Để hệ thống chạy liên tục (soak test). &
      Không treo, không rò bộ nhớ, bộ đếm gói tăng đều trong suốt thời gian thử.
      % TODO[cần đo]: thời lượng soak test + kết quả
      \\
    \bottomrule
  \end{tabular}
\end{table}

\subsection{Tổng hợp kết quả kiểm thử}
\label{section:4.4.5}

Bảng~\ref{tab:test-summary} tóm tắt trạng thái các hạng mục kiểm thử. Cột "Kết
quả" để tác giả điền sau khi thực hiện đầy đủ trên phần cứng.

\begin{table}[H]
  \centering
  \caption{Bảng tổng hợp kết quả kiểm thử.}
  \label{tab:test-summary}
  \begin{tabular}{@{}llcp{4.6cm}@{}}
    \toprule
    \textbf{Nhóm} & \textbf{Mã} & \textbf{Số kịch bản} & \textbf{Kết quả (Đạt/Không đạt)} \\
    \midrule
    Truyền thông CAN & TC\_CAN\_01--04 & 4 & % TODO[điền sau khi thử]
    \\
    Gateway CAN--Ethernet & TC\_ETH\_01--05 & 5 & % TODO
    \\
    Ứng dụng Android & TC\_APP\_01--05 & 5 & % TODO
    \\
    Tích hợp đầu cuối & TC\_E2E\_01--05 & 5 & % TODO
    \\
    \bottomrule
  \end{tabular}
\end{table}

\end{document}
